\documentclass[11pt,a4paper]{article}
\usepackage[margin=25mm]{geometry}
\usepackage{tikz}
\usepackage{booktabs}
\usepackage{longtable}
\usepackage{hyperref}
\usetikzlibrary{arrows.meta, positioning, shapes.geometric}
\title{Canonical CBOR codec --- recovery, validation, and publication plan}
\author{Codex integration plan}
\date{Created-at 2026-09-04 13:09 CEST}

\begin{document}
\maketitle

\section*{Authorisation}

\begin{quote}
\texttt{APPROVED-CORE-EDIT: canonical-cbor-codec}
\end{quote}

The first memory cell isolated a pre-existing per-incidence temporary in canonical hashing. The
output-preserving internal repair is separately authorised by:

\begin{quote}
\texttt{APPROVED-CORE-EDIT: canonical-hash-streaming-incidence}
\end{quote}

The user supplied this token verbatim on 2026-09-04 after an alternative Codex process had produced
an uncommitted implementation. The token authorises audit, correction, complete verification, and
publication of that exact core scope. The commit message repeats the token in a footer.

\section*{ETA}

Approximately 2.5 hours from approval: 30 minutes recovery and design audit, 45 minutes integration
and hostile-input tests, 35 minutes correctness and static gates, 25 minutes bounded memory/latency
evidence, and 15 minutes report, commit, and remote verification. The implementation already exists;
the uncertainty is validation depth rather than code volume.

\section*{1. Goal and ownership}

HyMeKo Core owns serialization and validation of its canonical compiled program. Transports such as
Chordex and Zenoh remain byte carriers. Consumers receive only a validated IR and reconstructed
interner:

\[
\texttt{CompiledProgram} \xrightarrow{\texttt{encode\_compiled}}
\texttt{CBOR bytes} \xrightarrow{\texttt{decode + limits}}
\texttt{DecodedCborPayload}.
\]

This work closes requirements \texttt{UIR-CBOR-01..03}. It extends the existing
\texttt{writers::cbor\_writer} owner and reuses \texttt{ir\_stream::DecodeLimits}; it does not add a
second codec, transport, IR, index, interner, or limit vocabulary.

\section*{2. Recovery context}

The implementation arrived while an older \texttt{origin/main} merge was unresolved. Before any
integration mutation, the working-only CBOR files, 724-line agent contract, corpus resolution, Git
index, and merge metadata were copied to
\texttt{D:/h/core-merge-safety-20260904-1253} and verified by SHA-256. The corpus resolution was
validated at 105 fixtures, staged alone, and the merge was closed as \texttt{d6be4ea}. CBOR work now
lives on \texttt{feat/canonical-cbor-codec}; \texttt{CLAUDE.md} remains outside its index.

\section*{3. Approved boundary}

\begin{longtable}{@{}lll@{}}
\toprule
Path & Change & Authority \\
\midrule
\texttt{hymeko\_core/src/writers/cbor\_writer.rs} & public facade and errors & approved core \\
\texttt{hymeko\_core/src/writers/cbor\_writer/*.rs} & preflight, validation, tests & approved core \\
\texttt{hymeko\_core/src/ir/canonical\_hash.rs} & bit-identical streaming incidence hash & approved core \\
\texttt{hymeko\_core/Cargo.toml} & direct pinned \texttt{ciborium 0.2.2} & approved core \\
\texttt{Cargo.lock} & dependency ownership update & approved core \\
\texttt{hymeko\_core/tests/**} & corpus/workload integration gates & allowlisted \\
\texttt{hymeko\_core/benches/**} & Criterion and peak-RSS evidence if needed & allowlisted \\
\texttt{docs/plans/**}, \texttt{reports/**} & traceability & allowlisted \\
\bottomrule
\end{longtable}

Out of scope: IR layout or canonical-hash schema/ordering changes; \texttt{ir\_stream} wire changes;
Chordex, Zenoh, simulator, or XR integration; mutation/snapshot APIs; a second serialization crate;
and the working-only \texttt{CLAUDE.md} migration.

The hash implementation change may only replace materialised incidence bytes with an equivalent
ordered stream. Frozen v1 known-answer vectors and the 105-fixture corpus digest are rollback gates:
any changed hash aborts the optimisation and the codec remains unpublished.

\section*{4. Contracts}

\paragraph{Encode.} Input paths must be non-empty, workspace-relative, Unicode, and portable across
supported hosts. The index, namespace paths, all IDs, strings, imports, hierarchy, and canonical hash
must be internally coherent. The output uses explicit schema version 1 and contains no trusted
derived state that decode accepts without reconstruction or comparison.

\paragraph{Decode.} Message, payload, text/byte strings, containers, and nesting are bounded before
Serde allocation. Indefinite-length items, truncation, trailing bytes, unsupported schema, invalid
portable paths, duplicate strings, invalid symbols/IDs, incoherent IR, unrebuildable or mismatched
index, namespace mismatch, and canonical-hash mismatch fail closed.

\paragraph{Postcondition.} A successful decode returns an immutable validated wrapper whose IR,
index, interner, imports, path, and canonical hash are mutually consistent. For valid corpus and
workload inputs, \texttt{decode(encode(x))} preserves canonical program identity.

\section*{5. Discriminating audit questions}

The inherited eight unit tests are necessary but not sufficient. Before accepting the code, answer
with executable tests:

\begin{enumerate}
  \item Are alias entries semantically validated, or merely copied after rebuilding \texttt{by\_path}?
  \item Are relative paths represented identically on Windows and Unix rather than serializing host
        separators into the wire contract?
  \item Does every truncated prefix and every indefinite container reject without panic?
  \item Do nested vectors and value trees obey pre-allocation limits rather than only top-level
        post-deserialization counts?
  \item Does corpus/workload round-trip preserve canonical hash, hierarchy, namespace roots, aliases,
        imports, and rendered semantics?
\end{enumerate}

A failed question is repaired minimally inside the approved owner and receives a regression test.

\section*{6. Test and evidence strategy}

\begin{itemize}
  \item Retain and extend module unit tests for schema, path, hash, index, ID, indefinite-length,
        trailing-byte, and truncation failures.
  \item Add one integration owner under the existing unified-IR validation suite; do not create a
        parallel test harness.
  \item Round-trip every lowerable corpus fixture and deterministic small/medium/large workload
        controls; compare semantic and canonical witnesses.
  \item Run exact CBOR tests, core library/tests, then full workspace tests.
  \item Run \texttt{cargo fmt --all -- --check} and pinned workspace Clippy
        \texttt{--all-targets -- -D warnings}.
  \item Measure Criterion encode/decode distributions on representative payload sizes only if no
        trustworthy existing cell covers the new facade. Record median, IQR, worst, source hash,
        fixture hash, and environment.
  \item Measure decode peak working set in a child process. Budget: less than twice decoded IR heap
        above process baseline and always below the 16 GiB host cap; no per-incidence temporary.
  \item Re-run the same baseline/materialise/decode child-process cells after the streaming hash
        repair; report both the failed pre-change cell and the accepted post-change cell.
\end{itemize}

\section*{7. Publication and rollback}

The merge commit and CBOR implementation remain separate. The CBOR commit stages only its approved
manifest, source, tests, plan, and report; the 724-line \texttt{CLAUDE.md} is excluded. Push
\texttt{main} only after confirming the merge head still descends from the unchanged remote main,
then push \texttt{feat/canonical-cbor-codec}. Verify local and remote object IDs.

Rollback is one revert of the CBOR commit. It removes the additive API/modules and direct dependency,
restoring the previous dormant payload definition. The already completed upstream merge is not part
of CBOR rollback.

\section*{8. Tooling waiver}

\texttt{tools.yaml} records that no TeX engine is installed on this Windows host. Therefore
\texttt{plan.tex}, \texttt{plan.tikz}, and \texttt{plan.mmd} are committed as renderable sources;
\texttt{plan.pdf} cannot be produced locally. This named waiver is repeated in the final report.

\section*{Architecture}

% Canonical CBOR codec ownership and trust boundary.
% Created-at 2026-09-04 13:09 CEST

\begin{tikzpicture}[
  font=\small,
  box/.style={draw, rounded corners=2pt, minimum height=9mm, align=center, inner sep=4pt},
  external/.style={box, fill=black!4},
  core/.style={box, fill=blue!7},
  validated/.style={box, fill=green!9},
  arr/.style={-Stealth, thick},
  reject/.style={box, fill=red!7},
  note/.style={font=\scriptsize\itshape, align=left, text width=40mm}
]

\node[core] (program) at (0,1.5) {\texttt{CompiledProgram} + \texttt{Interner}};
\node[core] (encode) at (4.2,1.5) {\texttt{encode\_compiled}\path normalization\witness verification};
\node[external] (bytes) at (8.3,1.5) {versioned CBOR bytes\transport-owned unchanged};
\draw[arr] (program) -- (encode);
\draw[arr] (encode) -- (bytes);

\draw[dashed, thick, red!60] (8.3,-3.0) -- (8.3,0.7)
  node[above, font=\scriptsize\bfseries, red!60] {untrusted input boundary};

\node[core] (preflight) at (11.8,0.5) {\texttt{preflight}\definite lengths\limits before Serde};
\node[core] (serde) at (11.8,-1.0) {Serde materialization\schema v1};
\node[core] (validate) at (7.4,-2.0) {semantic validation\IDs, paths, hierarchy, namespaces};
\node[core] (rebuild) at (3.2,-2.0) {rebuild index + interner\stream canonical hash};
\node[validated] (out) at (-1.1,-2.0) {\texttt{DecodedCborPayload}\validated IR + Interner};
\node[reject] (error) at (11.8,-2.8) {\texttt{CborError}\fail closed};

\draw[arr] (bytes) -- (preflight);
\draw[arr] (preflight) -- (serde);
\draw[arr] (serde) -- (validate);
\draw[arr] (validate) -- (rebuild);
\draw[arr] (rebuild) -- (out);
\draw[arr, red!60] (preflight) -- (error);
\draw[arr, red!60] (serde) -- (error);
\draw[arr, red!60] (validate) -- (error);

\node[note] at (3.2,-3.3) {Chordex, Zenoh, simulator, and XR are consumers or carriers; none owns this schema.};

\end{tikzpicture}


\end{document}
